\documentclass[10pt,a4paper]{amsart}
\usepackage[margin=19mm]{geometry}
\usepackage{amsmath,amssymb,mathtools}
\usepackage[scr=rsfs,cal=euler]{mathalfa}
\usepackage{xfrac}
\usepackage[unicode,hidelinks,bookmarks=false]{hyperref}
\newcommand{\ie}{i.e.\ }
\newcommand{\cf}{cf.\ }
\newcommand{\resp}{respectively\ }
\newcommand{\Subsection}{\subsection}
\providecommand{\nobreakdash}{\nobreak\hbox{-}}
\setlength{\emergencystretch}{2em}
\multlinegap0pt
\usepackage{bm}
\usepackage{bbm}
\usepackage{dsfont}
\usepackage{xspace}
\usepackage{cancel}
\usepackage{mathtools}
\usepackage{amsthm}
\makeatletter
\@ifundefined{theorem}{\newtheorem{theorem}{Theorem}}{}
\@ifundefined{lemma}{\newtheorem{lemma}[theorem]{Lemma}}{}
\@ifundefined{proposition}{\newtheorem{proposition}{Proposition}}{}
\makeatother
\newcommand{\N}{\mathbb{N}}
\title[On Landweber`s unique factorization problem]{On Landweber`s unique factorization problem}
\author{Adam Jones, Elad Paran}
\date{Source: arXiv:2607.03475v2 (CC BY 4.0)}
\begin{document}
\maketitle

\noindent\emph{Source and licence.} On Landweber`s unique factorization problem. Version arXiv:2607.03475v2 (CC BY 4.0), distributed under the Creative Commons Attribution 4.0 International licence, as recorded in the arXiv licence metadata for this exact version and shown on the abstract page https://arxiv.org/abs/2607.03475v2. Full rights notice: licenses/arxiv-2607.03475-CC-BY-4.0.txt.\par\smallskip

\noindent\textit{Editorial scope.} Counted unit: the Lemma whose source label is \texttt{None} in the article (source0.tex lines 543--560, source label \texttt{None}), delivered together with the article's own complete original proof of it (source0.tex lines 562--631, one \texttt{proof} environment of 70 lines). No mathematics is written, repaired or re-derived: statement and proof are the article's own text line for line. Required context retained uncounted: lem\_a|b (lines 187--187); quasi-prime (lines 455--457); theorem\_over\_dvr (lines 522--524). Every definition, display and label of those lines is retained unchanged. Omitted from this package and not redistributed: 1--186, 188--454, 458--521, 525--542, 561--561, 632--889. Nothing inside a retained excerpt is omitted or altered: every retained body is delivered exactly as the article's lines are written, so no omission receipt of any kind accompanies this package. Citations: the retained excerpts carry no citation command at all, so no bibliography entry of the article is transcribed and no reference is redistributed. Reference adaptation: none was needed and none was made; every reference inside the delivered unit resolves inside it, so no unresolved reference or citation is produced. Printed slips kept literal: no printed word, symbol, hypothesis or number of the counted statement or of its proof was altered. Layout and class: the article's own class is \texttt{article} with packages including inputenc, amsmath, amsthm, amsfonts, babel, mathtools, hyperref, comment, authblk, amscd, amssymb, bbm; the wrapper is the lane's pinned \texttt{amsart} editorial wrapper at 10pt on A4, which restates the theorem-like environments the retained text needs and adds the article's own macro definitions that the retained text needs (\texttt{\textbackslash N}), with extra wrapper packages (bm, bbm, dsfont, xspace, cancel, mathtools). The delivered numbering is the unit's own. Assets: no retained span contains a figure environment, a graphics-inclusion command, a PDF-page inclusion, a table or a TikZ picture; no image file is copied into this package, the package declares no asset dependency and no image file is redistributed with it. Licence: version arXiv:2607.03475v2 of this article is distributed under the Creative Commons Attribution 4.0 International licence, as recorded in the arXiv licence metadata for this exact version and shown on the abstract page https://arxiv.org/abs/2607.03475v2. A full rights notice is delivered at licenses/arxiv-2607.03475-CC-BY-4.0.txt. Characters: every retained excerpt is pure ASCII, so the delivered engine, XeLaTeX with its default Latin Modern text font, reports no missing character.
\begin{lemma}\label{lem_a|b} Let $R$ be a domain and let $a,b \in R[[t]]$. Then $a|_nb$ for all $n \in \N$ if and only if $a|b$ in $R[[t]]$. \end{lemma}

\begin{proposition}\label{quasi-prime}
For all $r,D\in\mathbb{N}$, there exists a constant $C_r(D)\in\mathbb{N}$ such that if $A\in V[[t]]$ is irreducible mod $t^D$, with $v(A_0)=r$, then $A$ is $C_r(D)$-prime. For $r=1$, we may take $C_r(D)=D$.
\end{proposition}

\begin{proposition}
    \label{theorem_over_dvr} Let $F \in V[[t]]$ be irreducible. Then $F$ is irreducible of finite height.
\end{proposition}

\begin{lemma}
Let $F_1,\ldots,F_s\in V[[t]]$  
be irreducible elements with non-zero constant terms. For each subset
$J\subseteq \{1,\ldots,s\}$, write
$$
F_J=\prod_{j\in J}F_j,
$$
with the empty product equal to $1$. Then, for every
$i\in\{0,\ldots,s\}$ and every $\ell \geq 1$, there exists an integer
$H_i(\ell)$ with the following property: if
$$
F_1\cdots F_i\equiv_n AB, n\ge H_i(\ell),
$$
then there is a subset $J\subseteq \{1,\ldots,i\}$ such that
$$
(A,t^\ell)=(F_J,t^\ell).
$$
\end{lemma}

\begin{proof}
We prove the statement by induction on $i$.

For $i=0$, take $H_0(\ell)=\ell$. Then $F_1\cdots F_i=1$, so 
$$
1\equiv_n AB, n\ge \ell,
$$
implies that $A$ is a unit modulo $t^\ell$ and hence
$(A,t^\ell)=V[[t]]$. 


Assume now that $i\ge 1$, and assume that the result has been proven for $i-1$. Put
$P=F_i,\ G=F_1\cdots F_{i-1}$ and let 
$
d_i=v(P_0)>0.
$
By Proposition \ref{theorem_over_dvr} there is $D_i\ge 1$ such that $P$ is irreducible mod $t^{D_i}$. Put $C_i=C_{d_i}(D_i)$
where $C_{d_i}(D_i)$ is the constant given by Proposition \ref{quasi-prime}. 

Choose $H_i(\ell)$ large enough so that for every $n\ge H_i(\ell)$, if
$$
n'=\left\lfloor \frac{n-C_i}{2}\right\rfloor,
$$
then
$$
n'\geq H_{i-1}(\ell),
n'\geq \ell.
$$


Suppose that $PG\equiv_n AB, n\ge H_i(\ell)$. 
Then $P|_n AB$. If $P|_r A$ for arbitrarily large $r$, then $P|A$ by Lemma \ref{lem_a|b}, and in particular $P|_{n'} A$. Thus we may assume that there exists a maximal integer $\alpha$ such that $P |_\alpha A$ . Similarly we may assume that $\beta$ is a maximal integer such that $P|_\beta B$. Then $$
\alpha+\beta>n-C_i.
$$
Therefore at least one of $\alpha,\beta$ is at least $n'$. Hence
$P |_{n'} A$ or $P|_{n'} B.$ First suppose that $P|_{n'}A$. Choose $A'\in V[[t]]$ such that $A \equiv_{n'} PA'$. Reducing $PG\equiv_n AB$ modulo $t^{n'}$, we get
$$
PG\equiv_{n'} PA'B.
$$
Since $P_0\neq 0$, we deduce that
$$
G\equiv_{n'} A'B.
$$
By the induction hypothesis applied to $G=F_1\cdots F_{i-1}$, there is a subset $J\subseteq \{1,\ldots,i-1\}$ such that
$$
(A',t^\ell)=(F_J,t^\ell).
$$

Since $n'\ge \ell$ and $A\equiv_{n'}PA'$, we have
$$
(A,t^\ell) =(PA',t^\ell) =
(PF_J,t^\ell) = (F_{J\cup\{i\}},t^\ell).
$$

Now suppose that $P|_{n'}B$. Choose $B'\in V[[t]]$ such that
$B\equiv_{n'}PB'$. Reducing $PG\equiv_n AB$ modulo $t^{n'}$, we get
$$
PG\equiv_{n'} APB'.
$$
and hence
$$
G\equiv_{n'} AB'.
$$
By the induction hypothesis applied to $G=F_1\cdots F_{i-1}$, there is a subset $J\subseteq \{1,\ldots,i-1\}$ such that
$$
(A,t^\ell)=(F_J,t^\ell).
$$

This proves the induction step and concludes the proof.
\end{proof}

\end{document}
